\scanpage{459}\mspage{371}
% archive sheet number: 458
b) \unsure{pour tout} $(I,\psi)\in\Lambda_2$, le foncteur composé
\diagram{\begin{tikzcd}[ampersand replacement=\&]
{}_\bullet I\arrow[r,"\psi"] \& C\arrow[r,"f"] \& \unsure{\mathcal N|N|\mathscr N}
\end{tikzcd}}
est $I$-exact à droite, i.e. fait de $(f\psi)(e_I)$ une limite projective de $(f\psi)|I$.

\deleted{\unsure{Sous}}\unsure{Avec} ces \deleted{conditions,} notations : si $\mathcal N$ est une catégorie accessible, il en est de même de $\underline{\mathrm{Hom}}_{\Lambda_1}^{\Lambda_2}(C,\mathcal N)$.

\uline{Commentaire} On utilisera ce théorème surtout quand on a $\Lambda_2=\varnothing$ ou $\Lambda_1=\varnothing$, auquel cas on note $\underline{\mathrm{Hom}}_{\Lambda_1}(C,\mathcal N)$, resp. $\underline{\mathrm{Hom}}^{\Lambda_2}(C,\mathcal N)$, la catégorie $\underline{\mathrm{Hom}}_{\Lambda_1}^{\Lambda_2}(C,\mathcal N)$. \deleted{De plus,} On appliquera le th. surtout quand $C$ elle-même est stable par $\unsure{\varinjlim|\underline{\lim}}$ de type $\Phi$ et \deleted{\unsure{par la}} $\varprojlim$ de type $\Psi$, et quand \deleted{les} $\Lambda_1$ (resp $\Lambda_2$) est déjà formé de foncteurs $\varphi$ (resp $\psi$) qui sont $I$-exacts à g. (resp. à droite). Donc $\underline{\mathrm{Hom}}_{\Lambda_1}^{\Lambda_2}(C,\mathcal N)$ est alors formée des $f:C\unsure{-|\longrightarrow}M$ qui commutent à certaines $\varinjlim$ de type $\Phi$, et à certaines $\varprojlim$ de type $\Psi$. En particulier, \unsure{quand}
